\chapter{SCF-DeepONet 模型优化与实验结果分析}
\label{ch:experiment_analysis}
\markboth{SCF-DeepONet 模型优化与实验结果分析}{}

在第三章中，本文完成了 SCF-DeepONet 的任务定义、数据表示、网络结构、多任务损失函数以及四阶段训练策略设计。为了验证所提出模型在道岔伤损 EMT 检测任务中的有效性，本章围绕实验设置、网络架构消融、监督基线对比、A 场辅助任务影响以及 PDE 弱物理微调效果展开分析。

本章首先介绍数据划分、训练配置及评价指标体系；随后通过与经典 Vanilla DeepONet 的消融对比，验证 SCF-DeepONet 架构设计对二维伤损检测和三维差场重构的贡献；在此基础上，构建纯热图监督与热图--场量联合监督两类基线模型，比较二者在同分布矩形伤损和圆形伤损零样本测试中的表现，从而分析三维磁矢势差场辅助监督对检测鲁棒性和分布外泛化能力的影响。最后，本章基于第三章定义的 PDE 弱物理残差，对 Stage 4 微调后的三维场重构结果进行配对统计分析，重点验证其对场预测物理一致性的改善作用，而不再重复 PDE 约束的数学推导。

\section{实验设置与评价指标}
\label{sec:experimental_setup}

为了客观评估 SCF-DeepONet 在二维伤损检测、三维场重构以及物理一致性方面的性能，本文从数据划分、训练设置、热图检测指标、三维场评价指标和配对统计分析方法五个方面构建实验体系。

\subsection{数据划分与训练设置}
\label{subsec:data_split}

本文所使用的矩形伤损数据集由程序控制有限元仿真软件批量生成，并按照物理样本级别进行划分。所谓物理样本级别划分，是指同一个伤损几何构型下的所有激励响应均只出现在同一个数据子集中，避免同一伤损样本在不同激励视角下同时进入训练集与验证集或测试集，从而造成信息泄漏。

在具体实验中，首先构建训练集与验证集，用于网络参数更新、阶段最优模型选择和超参数调整。其中，训练集包含 850 个物理伤损样本，验证集包含 150 个物理伤损样本。此外，本文单独保留一组在训练过程中完全不可见的独立测试集，共 65 个矩形伤损样本，用于最终性能评估。除矩形伤损测试外，本文还额外构建一组圆形伤损零样本测试集，共 50 个样本，用于评估模型对未见过伤损形状的分布外泛化能力。

所有模型均基于 PyTorch 框架实现，优化器采用 AdamW，并结合基于验证集指标的学习率调度机制控制训练过程\citep{loshchilov2017decoupled}。在纯监督主线训练中，基础学习率设置为 $3\times10^{-4}$；在后续热图精修或物理微调阶段，则进一步降低学习率至$1\times10^{-5}$，并仅对特定模块进行小幅参数更新。为保证不同实验组之间比较的公平性，本文所有对比实验均采用一致的数据划分、随机种子和阶段性模型选择规则。对于每个训练阶段，均根据该阶段的核心验证指标保存最优模型，而不是简单采用最后一个训练轮次的模型作为最终结果。

\subsection{热图检测评价指标}
\label{subsec:heatmap_metrics}

二维热图分支的任务是在探测平面上输出伤损区域的概率分布，因此本文采用交并比（Intersection over Union, IoU）作为热图检测性能的核心评价指标。设模型输出经过阈值化后的预测伤损区域为 $P$，真实伤损掩膜为 $G$，则 IoU 定义为
\begin{equation}
\label{eq:iou_definition}
IoU=\frac{|P\cap G|}{|P\cup G|}
\end{equation}

在 EMT 逆散射问题中，不同伤损尺寸、位置以及边界条件会导致样本难度呈现明显的长尾分布。若仅使用测试集平均 IoU 作为评价标准，容易掩盖模型在困难样本上的真实表现。因此，本文在热图检测任务中同时报告以下统计量：

\begin{itemize}
    \item \textbf{Mean IoU}：测试集中所有样本 IoU 的平均值，用于反映模型整体检测能力；
    \item \textbf{Median IoU}：测试集中所有样本 IoU 的中位数，用于反映典型样本上的检测水平；
    \item \textbf{Worst-10 Avg IoU}：测试集中 IoU 最低的 10 个样本的平均值，用于衡量模型在长尾困难工况下的保底检测能力。
\end{itemize}

对于圆形伤损零样本测试，由于其伤损几何形状与训练阶段所采用的矩形掩膜分布存在显著差异，因此上述指标还可用于评价模型在分布外形状条件下的迁移能力。

\subsection{三维场评价指标}
\label{subsec:field_metrics}

对于三维差场重构任务，本文采用均方根误差（Root Mean Square Error, RMSE）
衡量预测差场 $\Delta \hat{A}$ 与真实差场 $\Delta A$ 之间的一致性。
设单个样本中参与评估的三维点数为 $N_{3D}$，第 $j$ 个查询点为 $q_j$，
当前激励编号为 $e$，则三维差场重构 RMSE 定义为
\begin{equation}
\label{eq:rmse_A}
\mathrm{RMSE}_{A}
=
\sqrt{
\frac{1}{N_{3D}}
\sum_{j=1}^{N_{3D}}
\left\|
\Delta \hat{A}_{\mathrm{norm}}(q_j,e)
-
\Delta A_{\mathrm{norm}}(q_j,e)
\right\|_2^2
}
\end{equation}

其中，$\Delta \hat{A}_{\mathrm{norm}}(q_j,e)$ 和
$\Delta A_{\mathrm{norm}}(q_j,e)$ 均为 6 通道归一化磁矢势差场向量。

与二维检测任务类似，本文不仅报告测试集上的平均场误差（Mean RMSE），还同时统计中位数误差（Median RMSE）、90\% 分位数误差（P90 RMSE）以及误差最大的前 10 个样本均值（Worst-10 Avg RMSE）。其中，P90 RMSE 与 Worst-10 Avg RMSE 更适合刻画模型在复杂边界、极端样本或局部非物理畸变工况下的表现。

\subsection{配对统计分析方法}
\label{subsec:statistical_analysis}

在比较 NoPDE 与 PDE 模型时，PDE 弱物理约束对性能的影响通常表现为``小幅但一致''的变化。若仅比较测试集平均 RMSE，往往难以判断这种改善是否真正作用于多数样本。为此，本文引入逐样本配对统计分析方法。

具体而言，对同一测试样本分别计算 NoPDE 模型与 PDE 模型下的三维场重构误差，并定义误差差值为
\begin{equation}
\label{eq:delta_rmse}
\Delta RMSE
=
RMSE_{\mathrm{PDE}}
-
RMSE_{\mathrm{NoPDE}}
\end{equation}

当 $\Delta RMSE<0$ 时，说明引入 PDE 后该样本的场重构误差下降。通过统计误差下降样本数量占比，可以直观判断 PDE 是否在多数样本上产生正向作用。进一步地，本文采用 Wilcoxon 符号秩检验对配对误差分布进行非参数统计分析，以判断 PDE 引入后的性能改善是否具有统计学显著性\citep{wilcoxon1945individual}。该方法不要求误差差值服从高斯分布，适合本文这类样本数量相对有限、误差分布可能存在长尾的实验场景。

\section{网络架构对比实验}
\label{sec:ablation_experiment}

在分析 SCF-DeepONet 训练策略与辅助任务作用之前，首先需要回答一个更根本的问题：第三章所提出的网络架构设计——包括共享条件化编码、Fourier 位置编码、FiLM 条件调制、CNN 特征图精修以及热图--场耦合机制——究竟为 EMT 伤损检测带来了多少实质性提升。本节通过与经典 DeepONet 基线的消融对比，对上述问题展开分析。

\subsection{消融实验设计}
\label{subsec:ablation_design}

为验证 SCF-DeepONet 架构设计的有效性，本节构建两个对比模型：

\textbf{（1）SCF-DeepONet（本文方法）。}  
即第三章所提出的完整网络架构，包含共享观测编码模块（PairEncoder + Pool + Refiner）、Fourier 位置编码、FiLM 条件调制、CNN 特征图精修、注意力池化以及二维--三维热图耦合机制。

\textbf{（2）VanillaDeepONet（经典基线）。}  
该模型严格遵循 Lu 等人\upcite{lu2021learning} 提出的原始 Branch--Trunk 点积结构，仅作面向 EMT 双任务的最小必要适配：Branch 采用 Flatten + MLP 编码输入观测为单一潜在特征向量 $z$；二维热图 Trunk 和三维场 Trunk 分别通过坐标输入与 $z$ 的点积运算输出预测结果。该基线不包含任何本文提出的架构创新组件（多层级 latent、傅里叶编码、FiLM、CNN 精修、热图--场耦合等）。

为保证对比的公平性，两个模型在消融实验中采用完全一致的实验条件：相同的数据划分（850/150/65）、相同的四阶段训练流程、相同的 Stage 3 评估协议（冻结共享编码与场分支，仅精修热图输出头）、相同的学习率策略以及相同的评价指标。由此，Stage 3 上的性能差异可单纯归因于网络架构本身的不同。

\subsection{Stage 3 热图检测与三维场重构性能对比}
\label{subsec:ablation_stage3_results}

表 \ref{tab:ablation_comparison} 给出了 SCF-DeepONet 与 VanillaDeepONet 在独立测试集上的 Stage 3 性能对比，涵盖热图检测指标（IoU）与三维场重构指标（RMSE）两个维度。

\begin{table}[htbp]
    \centering
    \tabcaption{SCF-DeepONet 与 VanillaDeepONet Stage 3 性能对比}{Stage 3 Performance Comparison: SCF-DeepONet vs VanillaDeepONet}
    \label{tab:ablation_comparison}

    \renewcommand{\arraystretch}{1.25}
    \zihao{5}

    \begin{threeparttable}
    \begin{tabularx}{0.98\textwidth}{
        >{\centering\arraybackslash}p{0.12\textwidth}
        >{\centering\arraybackslash}p{0.09\textwidth}
        >{\centering\arraybackslash}p{0.09\textwidth}
        >{\centering\arraybackslash}p{0.09\textwidth}
        >{\centering\arraybackslash}p{0.09\textwidth}
        >{\centering\arraybackslash}p{0.09\textwidth}
        >{\centering\arraybackslash}p{0.09\textwidth}
        >{\centering\arraybackslash}p{0.09\textwidth}
    }
        \toprule
        模型 & Mean IoU & Median IoU & Worst-10 Avg IoU & Mean RMSE & Median RMSE & P90 RMSE & Worst-10 Avg RMSE \\
        \midrule
        Vanilla DeepONet & 67.79\% & 73.52\% & 37.09\% & 1.82e-1 & 1.80e-1 & 2.24e-1 & 2.32e-1 \\
        SCF-DeepONet & 81.53\% & 85.57\% & 51.34\% & 1.77e-1 & 1.72e-1 & 2.17e-1 & 2.23e-1 \\
        \bottomrule
    \end{tabularx}
    \begin{tablenotes}
        \footnotesize
        \item 注：热图检测指标为 IoU，三维场重构指标为 RMSE（归一化量纲）。
    \end{tablenotes}
    \end{threeparttable}
\end{table}

从热图检测指标来看，在完全一致的 Stage 3 精修条件下，SCF-DeepONet 的各项 IoU 指标均显著优于 VanillaDeepONet。具体而言，SCF-DeepONet 的 Mean IoU 相较 VanillaDeepONet 提升了约 13.73\%，Median IoU 提升了约 12.04\%，而 Worst-10 Avg IoU 提升了约 14.24\%。Worst-10 指标的改善尤为关键，它说明 SCF-DeepONet 的架构设计不仅提升了典型样本上的检测精度，更增强了对复杂边界和长尾困难工况的鲁棒性。

从三维场重构指标来看，SCF-DeepONet 在 Mean RMSE、Median RMSE、P90 RMSE 以及 Worst-10 Avg RMSE 上同样全面低于 VanillaDeepONet。这一结果表明，第三章所提出的多层级条件编码、Fourier 位置编码以及 FiLM 调制等组件，并非冗余设计，而是针对 EMT 多激励观测与连续场回归问题特点的必要架构扩展——它们不仅改善了二维热图检测，也同步提升了三维差场重构的质量。

进一步，图 \ref{fig:ablation_cdf} 给出了两模型在测试集上 IoU 的累积分布函数。SCF-DeepONet 的曲线在所有 IoU 阈值上均位于 VanillaDeepONet 右方，说明其性能优势不仅体现在均值层面，更覆盖了整个样本分布——包括困难样本尾部。图 \ref{fig:ablation_heatmap_grid} 展示了代表性样本的热图预测可视化对比，可以直观观察到 SCF-DeepONet 的预测边界更加清晰、与真实伤损框的重合度更高，而 VanillaDeepONet 的预测存在更为明显的边界不贴合及伪影现象。

\begin{figure}[htbp]
    \centering
    \includegraphics[width=0.76\textwidth]{figures/fig_ablation_cdf.png}
    \figcaption{SCF-DeepONet 与 VanillaDeepONet 的 IoU 累积分布对比}{CDF Comparison of IoU between SCF-DeepONet and VanillaDeepONet}
    \label{fig:ablation_cdf}
\end{figure}

\begin{figure}[ht]
    \centering
    \includegraphics[width=0.78\textwidth]{figures/fig_ablation_heatmap_grid.jpg}
    \figcaption{代表性样本热图预测对比}{Representative Heatmap Prediction Comparison}
    \label{fig:ablation_heatmap_grid}
\end{figure}

\subsection{消融实验结论}
\label{subsec:ablation_conclusion}

综合 Stage 3 消融实验的定量和定性结果，可以得出以下结论：

\textbf{（1）本文所提出的 SCF-DeepONet 架构设计是改善铁磁性 EMT 伤损检测性能的关键因素。} 在完全相同的训练流程、数据集和评估协议下，SCF-DeepONet 相较经典 Vanilla DeepONet 基线在热图检测 IoU 上取得了明显提升（Mean IoU 提升约 13.73\%，Worst-10 Avg IoU 提升约 14.24\%）。这一提升幅度明确表明，本文针对 EMT 多激励结构化输入和连续域输出所设计的共享条件编码、Fourier 位置编码、FiLM 调制、CNN 特征图精修以及热图--场耦合机制，是适配铁磁性 EMT 逆问题特点的必要架构扩展，而非冗余设计。

\textbf{（2）性能优势覆盖全样本分布。} SCF-DeepONet 不仅在 Mean IoU 上大幅领先，在 Worst-10 Avg IoU 以及 CDF 尾部同样保持显著优势；三维场重构指标在 P90 和 Worst-10 尾部也全面优于基线。证明其架构设计增强了模型对复杂边界和长尾工况的鲁棒性，而非仅提升了``容易样本''的拟合能力。

\textbf{（3）后续训练策略分析建立在可靠架构基线之上。} 在确认网络架构本身具有显著优越性的前提下，下节中对训练策略（Stage 演进、A 场辅助任务作用）的分析将建立在已验证的架构基线之上，避免了``训练策略的收益可能来源于架构本身''的混淆。换言之，后文所讨论的辅助监督与 PDE 约束，是在已经具备较好检测能力的架构基础上进一步探索的增强手段，而非弥补架构缺陷的补救措施。

\section{监督基线构建}
\label{sec:supervised_baselines}

本节构建纯热图监督与热图--场量联合监督两类基线模型，作为后续 A 场辅助任务分析、分布外泛化实验以及 PDE 物理约束实验的对比参照。

\textbf{纯热图监督基线：} 仅使用二维伤损热图掩膜而弃用三维差场真值，经 Stage 1 预训练与 Stage 3 精修得到。该基线代表了在不借助三维场辅助信息的条件下，模型在``电压响应 $\rightarrow$ 二维伤损热图''映射任务上的检测能力。

\textbf{热图--场量联合监督基线：} 在 Stage 2 中引入三维差场辅助监督，并同样经 Stage 3 精修得到。与纯热图基线相比，该模型通过多任务联合训练使共享特征空间保留了与三维电磁散射结构相关的判别性信息。

两类基线采用一致的 SCF-DeepONet 架构、相同的数据划分和精修协议，以保证对比的公平性。各 Stage 的训练目标与参数冻结策略已在第 \ref{subsec:stage_objectives} 节中给出，本文后续实验均基于该训练策略展开，此处不再重复说明。

\section{A 场辅助任务对伤损检测性能的影响}
\label{sec:impact_of_a_field}

在确立纯热图监督模型与热图--场量联合监督模型之后，本节进一步通过消融实验分析三维差场辅助任务对二维伤损检测性能的实际影响。需要强调的是，本节关注的重点并非三维场本身的重构精度，而是三维差场辅助监督是否改变了模型对二维伤损检测任务的特征提取方式，以及这种改变是否能够提升模型在困难样本和分布外形状下的鲁棒性。

\subsection{同分布矩形伤损上的对比}
\label{subsec:id_comparison}

首先，在与训练分布一致的矩形伤损测试集上对两类模型进行比较。表 \ref{tab:heatmap_id_ood_comparison} 给出了两类模型在同分布矩形测试集和圆形伤损零样本测试集上的热图评价指标。

从同分布矩形测试结果可以看出，纯热图监督模型在 Mean IoU 和 Median IoU 上分别达到 $75.16\%$ 和 $80.73\%$，略高于热图--场量联合监督模型的 $73.82\%$ 和 $78.25\%$。这说明，在训练分布内部，纯热图模型由于所有参数更新均直接服务于二维检测任务，能够形成较高效的同分布拟合能力。A 场辅助任务在此场景下并未以牺牲 ID 精度为代价换取泛化，而是体现了其作为\textbf{训练阶段正则化机制}的定位：它并不追求在熟悉样本上更激进地拟合，而是通过三维场约束迫使共享特征保留更具物理泛化性的表征，这也是后续 PDE 物理约束能够生效的前提。

\begin{table}[htbp]
    \centering
    \tabcaption{不同测试集模型性能指标对比}{Comparison of model performance metrics on different test sets}
    \label{tab:heatmap_id_ood_comparison}

    \renewcommand{\arraystretch}{1.25}
    \zihao{5}

    \begin{threeparttable}
    \begin{tabularx}{0.98\textwidth}{
        >{\centering\arraybackslash}p{0.12\textwidth}
        >{\centering\arraybackslash}p{0.10\textwidth}
        >{\centering\arraybackslash}p{0.10\textwidth}
        >{\centering\arraybackslash}p{0.10\textwidth}
        >{\centering\arraybackslash}p{0.12\textwidth}
        >{\centering\arraybackslash}p{0.10\textwidth}
        >{\centering\arraybackslash}p{0.10\textwidth}
        >{\centering\arraybackslash}p{0.10\textwidth}
    }
        \toprule
        \multirow{2}{*}{模型} &
        \multicolumn{3}{c}{矩形伤损测试集（ID）} &
        \multicolumn{3}{c}{圆形伤损测试集（OOD）} \\
        \cmidrule(lr){2-4}\cmidrule(lr){5-7}
        & Mean IoU & Median IoU & Worst-10 Avg IoU
        & Mean IoU & Median IoU & Worst-10 Avg IoU \\
        \midrule
        纯热图
        & $75.16\%$ & $80.73\%$ & $42.90\%$
        & $34.05\%$ & $40.47\%$ & $0.00\%$ \\
        热图--场量联合
        & $73.82\%$ & $78.25\%$ & $46.16\%$
        & $40.31\%$ & $46.02\%$ & $2.33\%$ \\
        \bottomrule
    \end{tabularx}
    \end{threeparttable}
\end{table}

图 \ref{fig:rectangular_samples} 给出了矩形测试集中代表性样本的可视化结果。图中每一列对应同一个样本，第一行为真实伤损位置，二三行分别为两类模型的预测热图。可以看出，在表现较好的样本上，两类模型均能够较准确地覆盖真实伤损掩膜；但在困难样本中，纯热图监督模型更容易出现局部偏移、伪热点或边界响应分散等现象，而热图--场量联合监督模型的预测区域相对更加集中，位置也更接近真实伤损区域。这一现象与表 \ref{tab:heatmap_id_ood_comparison} 中 Worst-10 Avg IoU 的统计结果一致。

\begin{figure}[ht]
    \centering
    \includegraphics[width=0.78\textwidth]{figures/TestSet_Ranks.jpg}
    \figcaption{矩形测试集代表性样本预测结果}{Representative prediction results on the rectangular test set}
    \label{fig:rectangular_samples}
\end{figure}

为了进一步从传统 EMT 成像方法与深度学习方法的角度进行对比，图 \ref{fig:traditional_comparison} 给出了包括LBP、Tikhonov、Landweber三种EMT经典反演算法与 SCF-DeepONet 在十个随机测试样本上的重建结果。其中每行为一个样本，每列是不同的算法效果，第一列为真实伤损掩膜。

\begin{figure}[p]
    \centering
    \includegraphics[width=\textwidth,height=0.9\textheight,keepaspectratio]{figures/Traditional_vs_DL_10Samples_Transposed.jpg}
    \figcaption{传统算法与SCF-DeepONet检测结果对比}{Comparison of traditional and SCF-DeepONet detection results}
    \label{fig:traditional_comparison}
\end{figure}

从图 \ref{fig:traditional_comparison} 可以看出，传统线性重建方法能够在一定程度上反映伤损扰动趋势，但其结果通常存在较明显的扩散、伪影和边界模糊现象；相比之下，SCF-DeepONet 直接从多视角电压响应中学习非线性映射关系，能够输出更加集中的伤损概率区域。该对比说明，在铁磁材料 EMT 这种强非线性、病态性较强的问题中，深度学习模型具有更强的特征整合能力和检测表达能力。

\subsection{圆形伤损零样本测试中的 OOD 泛化对比}
\label{subsec:ood_comparison}

为了进一步检验模型是否真正学习到了具有物理意义的共享表征，本文构建了圆形伤损零样本测试集。该测试集中的伤损几何形状与训练阶段使用的矩形掩膜明显不同，因此能够有效考察模型面对未见过几何边界时的迁移能力。

由表 \ref{tab:heatmap_id_ood_comparison} 可以看出，面对圆形伤损 OOD 测试，两类模型的 IoU 均较同分布矩形测试明显下降，说明几何分布偏移确实给模型检测带来了显著挑战。但相较于纯热图监督模型，热图--场量联合监督模型在三项指标上均取得更好表现：其 Mean IoU 为 $40.31\%$，高于纯热图监督模型的 $34.05\%$；Median IoU 为 $46.02\%$，高于纯热图监督模型的 $40.47\%$；Worst-10 Avg IoU 为 $2.33\%$，而纯热图监督模型在该指标上为 $0.00\%$。

图 \ref{fig:circular_ood_samples} 给出了圆形伤损零样本测试中的代表性预测结果。可以看出，两个模型在缺陷尺寸较小且靠近边缘的困难样本中均出现了明显的响应漂移和伪热点现象，预测区域与真实圆形伤损位置之间存在较大偏差。但相比之下，热图--场量联合监督模型对圆形缺陷大小的预测更加准确，说明三维差场辅助任务小幅增强了模型对未知形状的适应能力。另外，LBP、Tikhonov、Landweber等EMT经典图像重建算法的表现与矩形集上几乎一致，故不再给出单独对比图。

\begin{figure}[htbp]
    \centering
    \includegraphics[width=0.80\textwidth]{figures/ood_ranks.jpg}
    \figcaption{圆形伤损零样本测试集代表性样本预测结果}{Representative prediction results on the circular OOD test set}
    \label{fig:circular_ood_samples}
\end{figure}

这一结果具有重要意义。若模型仅依赖二维矩形热图掩膜进行训练，则其特征提取过程更容易围绕``矩形几何模板''形成统计相关性；一旦测试样本的真实形状发生改变，模型原先建立的映射关系便容易失效。相反，当引入三维差场辅助任务后，共享观测编码模块必须同时解释二维几何伤损与三维电磁散射场之间的对应关系，从而被迫学习更加底层、更加具有物理约束性的共享表征。这种表征虽然不一定在训练分布内带来更高的平均精度，却能在分布外场景中表现出更好的迁移能力。

\subsection{A 场辅助任务的作用机制分析}
\label{subsec:mechanism_analysis}

综合同分布矩形测试与圆形伤损零样本测试结果，可以更清晰地理解三维 A 场辅助任务在本文框架中的实际作用。

首先，A 场预测并非系统部署阶段必须直接输出的最终工程结果，其更核心的价值体现在训练阶段。通过在主检测任务之外增加三维差场回归约束，网络前端的共享观测编码模块不再只是学习``多通道观测响应到二维热图''的直接映射，而是必须进一步保留与三维电磁散射结构相关的信息。因此，A 场辅助任务本质上是一种\textbf{训练阶段的中间监督与特征正则化机制}。

其次，从实验结果来看，A 场辅助任务\textbf{并未}在同分布矩形样本上取得更高的平均检测精度（Mean IoU 与 Median IoU 略低于纯热图模型），但在长尾困难样本（Worst-10 Avg IoU 提升约 3.26\%）和圆形 OOD 样本（Mean IoU 从 34.05\% 提升至 40.31\%）上表现出了更稳定的检测下限和更好的泛化能力。这说明，A 场辅助任务的主要收益并不体现在``让模型在熟悉样本上更激进地拟合''，而体现在\textbf{抑制模型对训练形状模板的过拟合，并提升未知伤损形状的鲁棒性}。这一结论也与第三章中将 A 场任务定位为``训练阶段辅助机制''的设计意图一致。

需要指出的是，A 场辅助任务的引入也带来了额外训练代价。三维差场真值包含多激励、多空间点和复数场分量，其数据规模明显大于二维热图掩膜；同时，Stage 2 多任务训练需要额外的场量回归过程，且场真值依赖有限元仿真。因此，A 场辅助监督更适合作为训练阶段的物理感知正则化机制，而不是系统部署阶段必须输出的工程检测结果。

此外，A 场辅助监督与后续 PDE 弱物理微调的作用层次不同。前者通过有限元差场真值约束共享特征，使模型在检测任务中获得更好的困难样本鲁棒性和分布外泛化能力；后者则基于第三章定义的 PDE 相对残差，仅在 Stage 4 中对三维场分支进行后期物理一致性修正。就本文实验结果而言，二维检测性能的变化主要来自 Stage 2 的多任务联合训练，而 PDE 微调并不直接改变二维热图分支。因此，A 场辅助监督服务于检测泛化性能，PDE 弱约束服务于场预测的物理一致性，二者目标不同但在训练流程上具有递进关系。

\section{PDE 弱物理微调实验与物理一致性分析}
\label{sec:pde_implementation}

前文实验表明，三维差场辅助监督能够改善模型在困难样本和分布外形状下的检测鲁棒性。然而，Stage 1 至 Stage 3 仍主要依赖有限元仿真真值进行监督训练，三维场预测结果虽然在数值上接近仿真差场，但并未显式约束其对频域电磁控制方程的满足程度。为进一步考察物理约束对三维场分支的修正作用，本文在 Stage 4 中基于第 \ref{subsec:pde_loss} 节定义的 PDE 相对残差损失，对三维场分支进行弱物理微调实验。

需要强调的是，本节关注的对象是三维磁矢势差场重构结果，而不是二维伤损热图检测结果。按照第 \ref{subsec:stage_objectives} 节给出的训练策略，Stage 4 冻结共享观测编码模块与二维热图分支，仅微调 FieldTrunk。因此，PDE 弱物理约束不会直接改变二维热图预测过程，其主要目标是提升三维场预测的物理一致性和局部稳定性。

\subsection{PDE 微调实验设置}
\label{subsec:pde_finetune_setting}

为验证 PDE 弱物理约束的实际作用，本文以 Stage 3 训练完成后的热图--场量联合监督模型作为初始化模型，并在 Stage 4 中继续微调三维场分支。PDE 残差计算所需的总场构造、软伤损掩膜、材料参数插值以及相对残差归一化方式均沿用第 \ref{subsec:pde_loss} 节的定义。微调过程中，网络预测的归一化差场首先与无伤损基准场叠加得到总场，再恢复至真实物理量纲，并在源-free 区域采样点上计算 PDE 相对残差。

实验中设置两类对比模型：

\textbf{（1）NoPDE 模型。}  
该模型为 Stage 3 结束后的热图--场量联合监督模型，不再引入 PDE 物理残差。它代表仅依靠有限元差场真值监督时的三维场重构水平。

\textbf{（2）PDE 模型。}  
该模型以 NoPDE 模型为初始化，在 Stage 4 中加入 PDE 相对残差损失，并仅更新 FieldTrunk 参数。由于二维热图分支和共享编码模块被冻结，该模型与 NoPDE 模型在二维检测路径上保持一致，差异主要来自三维场分支的物理微调。

为避免平均指标掩盖逐样本差异，本文采用第 \ref{subsec:statistical_analysis} 节所述的配对统计方法，对 65 个独立矩形测试样本分别计算 NoPDE 与 PDE 模型的三维场重构 RMSE，并统计逐样本误差变化情况。

\subsection{三维场 RMSE 配对结果}
\label{subsec:pde_rmse_results}

图 \ref{fig:pde_scatter} 给出了 NoPDE 与 PDE 模型的逐样本 RMSE 配对散点图。图中虚线表示 $y=x$，即两种模型在同一样本上的三维场重构误差相同。若样本点位于虚线下方，则表示引入 PDE 弱物理约束后该样本的 RMSE 下降。

\begin{figure}[!hthp]
    \centering
    \includegraphics[width=0.68\textwidth]{figures/Scatter_NoPDE_vs_PDE.jpg}
    \figcaption{NoPDE 与 PDE 模型逐样本 RMSE 配对散点图}{Scatter plot of paired RMSE for NoPDE vs PDE models}
    \label{fig:pde_scatter}
\end{figure}

可以看出，绝大多数样本点位于参考线下方，说明 PDE 微调在多数测试样本上均降低了三维场重构误差。统计结果显示，在 65 个独立测试样本中，有 62 个样本在引入 PDE 弱约束后 RMSE 出现下降，占比约为 95.4\%。进一步采用 Wilcoxon 符号秩检验对逐样本误差差值进行非参数统计分析，结果表明该配对差异具有统计学显著性（$p<0.05$）。这说明 PDE 弱约束并非只在个别样本上产生偶然改善，而是在多数样本上形成了稳定一致的场误差压制效果。

从误差改善幅度看，PDE 微调带来的 RMSE 下降整体较为温和。这一现象与 Stage 4 的训练定位一致：PDE 弱约束并不是重新学习完整的电压响应到三维场分布映射，而是在已有监督训练结果基础上，对三维场分支进行保守式后期修正。因此，其作用更接近于削弱局部非物理波动、改善场量预测的一致性，而不是大幅改变三维场整体形态。

\subsection{PDE 约束作用与局限性分析}
\label{subsec:pde_effect_limitation}

PDE 微调实验表明，第三章所定义的相对 PDE 残差能够在铁磁性 EMT 场景下稳定参与训练，并对三维场预测产生小幅但一致的改善。这说明在已有 A 场辅助监督初始化的基础上，引入弱物理残差是一条可行的物理一致性增强路径。

不过，也需要明确其作用边界。首先，Stage 4 中仅 FieldTrunk 参与参数更新，共享观测编码模块和二维热图分支保持冻结，因此 PDE 微调不会直接改善二维伤损检测 IoU。实际结果也表明，引入 PDE 后二维热图检测指标与 Stage 3 精修结果基本一致，未出现显著变化。其次，本文采用的是弱物理微调策略，PDE 损失权重较小，且仅作为后期正则项参与优化，因此其主要贡献体现在三维场预测的物理一致性层面，而非检测精度的进一步提升。

综上，PDE 弱物理微调在本文框架中的作用可概括为：在 SCF-DeepONet 已经具备较好二维检测能力和三维差场重构能力的基础上，利用频域电磁控制方程残差对三维场分支进行后期修正，使场预测结果在多数样本上获得统计显著的误差下降。该阶段不是独立的检测性能来源，而是对模型物理可靠性和可解释性的补充增强。

\section{本章小结}
\label{sec:ch4_conclusion}

本章围绕 SCF-DeepONet 模型的训练优化与实验验证展开，从网络架构消融、监督基线构建、A 场辅助任务分析和 PDE 弱物理微调四个方面进行了系统评估。

首先，通过与经典 Vanilla DeepONet 的消融对比，验证了本文所提出的 SCF-DeepONet 架构设计在铁磁性 EMT 伤损检测中的核心价值。在完全相同的训练条件下，SCF-DeepONet 的 Mean IoU 和 Worst-10 Avg IoU 均明显优于 Vanilla DeepONet，且性能优势覆盖整个样本分布。这一结果表明，共享条件编码、Fourier 位置编码、FiLM 调制、CNN 特征图精修以及热图--场耦合机制，是适配 EMT 多激励结构化输入和连续域输出需求的关键架构设计，也是本文提升检测性能的主要来源。

在此基础上，本章进一步分析了三维差场辅助监督对二维伤损检测的实际影响。实验表明，A 场辅助任务虽未在同分布矩形样本上取得更高的平均 IoU，但在困难样本检测下限和圆形伤损零样本泛化能力上表现出更稳定的优势。这说明 A 场辅助监督的核心价值不在于直接提升同分布检测精度，而在于抑制模型对训练形状模板的过拟合，并提升未知伤损形状的鲁棒性。同时，该阶段也为后续 PDE 弱物理微调提供了较好的三维场初始化基础。

最后，本章讨论了 PDE 弱物理约束在三维场分支上的应用效果。实验表明，PDE 微调能够在不改变二维检测路径的前提下，使多数测试样本的三维场重构误差获得小幅而统计显著的下降。由于 Stage 4 仅微调 FieldTrunk，其贡献主要体现在场预测的物理一致性层面，而不是二维检测精度的进一步提升。

综上所述，本章明确了本文各技术组件的贡献层次：SCF-DeepONet 架构设计是提升铁磁性 EMT 伤损检测性能的核心基础；三维差场辅助监督是增强泛化鲁棒性的训练机制，并为后续物理约束嵌入提供初始化条件；PDE 弱物理微调则是在此基础上对三维场预测物理一致性的补充修正。三者构成了``主架构设计 + 辅助场监督 + 后期物理微调''的递进式智能反演方案。